\documentclass[10pt]{article}
\usepackage[T1]{fontenc}
\usepackage[utf8]{inputenc}
\usepackage[french]{babel}
\usepackage{amsmath,amssymb,mathrsfs,mathpazo}
\usepackage{geometry}
\geometry{a4paper,margin=2cm}
\usepackage{graphicx,array,tabularx,tabularray,colortbl,multirow,multicol,booktabs,enumitem,pifont,ulem,fancybox,diagbox,cancel}
\usepackage{pstricks,pstricks-add,pst-plot,pst-tree,pst-node,pst-text,pst-func,pst-math,pst-bspline,pst-3dplot,pst-eucl,pst-solides3d}
\let\pscalculate\undefined
\usepackage{pst-bezier}
\usepackage{tikz,pgfplots,tkz-tab}
\usetikzlibrary{arrows,arrows.meta,patterns,calc,positioning,trees,decorations.markings,shapes.geometric,shapes.misc,shapes.arrows,intersections,decorations.pathmorphing,decorations.pathreplacing,backgrounds,shadows,babel}
\pgfplotsset{compat=1.18}
\usepgfplotslibrary{fillbetween}
\usepackage[np]{numprint}
\usepackage{listings,makecell,pas-tableur,twoopt}
\let\barre\undefined
\usepackage[pstricks]{tablvar}
\definecolor{darkgreen}{RGB}{0,100,0}
\providecommand{\red}{\color{red}}
\providecommand{\blue}{\color{blue}}
\providecommand{\green}{\color{green}}
\providecommand{\black}{\color{black}}
\providecommand{\white}{\color{white}}
\providecommand{\gray}{\color{gray}}
\providecommand{\cyan}{\color{cyan}}
\providecommand{\magenta}{\color{magenta}}
\providecommand{\yellow}{\color{yellow}}
\providecommand{\violet}{\color{violet}}
\providecommand{\orange}{\color{orange}}
\providecommand{\euro}{\texteuro}
\providecommand{\EUR}{\texteuro}
\providecommand{\ieme}{\textsuperscript{e}}
\providecommand{\degres}{\ensuremath{^\circ}}
\providecommand{\degree}{\ensuremath{^\circ}}
\providecommand{\vv}[1]{\overrightarrow{#1}}
\providecommand{\Coord}[1]{#1}
\providecommand{\decosix}{}
\providecommand{\starredbullet}{\ensuremath{\bullet}}
\pagestyle{empty}
\setlength{\parindent}{0pt}
\usepackage{adjustbox}\DeclareUnicodeCharacter{2236}{:}\DeclareUnicodeCharacter{2190}{\ensuremath{\leftarrow}}
\providecommand{\R}{\mathbb{R}}
\providecommand{\N}{\mathbb{N}}
\providecommand{\Z}{\mathbb{Z}}
\providecommand{\Q}{\mathbb{Q}}
\providecommand{\D}{\mathbb{D}}
\providecommand{\C}{\mathbb{C}}
\providecommand{\e}{{\mathrm{e}}}
\providecommand{\vect}[1]{\overrightarrow{#1}}
\providecommand{\vectt}[1]{\overrightarrow{\mathrm{#1}}}
\providecommand{\vvt}[1]{\overrightarrow{\mathrm{#1}}}
\providecommand{\np}[2][]{{#2~#1}}
\providecommand{\vv}[1]{{\overrightarrow{#1}}}
\renewcommand{\d}{{\mathrm{d}}}
\providecommand{\degres}{{^\circ}}
\providecommand{\degre}{{^\circ}}
\providecommand{\blue}{\color{blue}}
\providecommand{\red}{\color{red}}
\providecommand{\enskip}{\quad}
\providecommand{\llbracket}{\lbrack\!\lbrack}
\providecommand{\rrbracket}{\rbrack\!\rbrack}
\providecommand{\euro}{{\text{€}}}
\providecommand{\no}{{\text{n°}}}
\providecommand{\interval}[2]{{[#1;#2]}}
\providecommand{\intervFO}[2]{{[#1;#2[}}
\providecommand{\dotfill}{\ldots}

\providecommand{\ds}{\displaystyle}
\providecommand{\textding}[1]{#1}
\providecommand{\Oij}{\(\mathrm{O};\vec{\imath},\vec{\jmath}\)}
\providecommand{\Oijk}{\(\mathrm{O};\vec{\imath},\vec{\jmath},\vec{k}\)}
\providecommand{\Ouv}{\(\mathrm{O};\vec{u},\vec{v}\)}
\providecommand{\intervFF}[2]{\ensuremath{[#1;#2]}}
\renewcommand{\euro}{\ifmmode\text{\texteuro}\else\texteuro\fi}
\DeclareUnicodeCharacter{00B0}{\ifmmode^{\circ}\else\textdegree\fi}
\DeclareUnicodeCharacter{20AC}{\ifmmode\text{\texteuro}\else\texteuro\fi}
% Énoncés : Éducation nationale. Transcription et corrigés : APMEP, Denis Vergès et contributeurs. Découpage et classement : CoopMaths, Sébastien Lozano, Éric Elter et contributeurs. Adaptation de la présentation : Maths Explorer.
% CC BY-SA : https://coopmaths.fr/www/about/
% https://coopmaths.fr/alea/static/bac/2024/tex/bac_2024_01_sujet0_1.tex
\begin{document}
\noindent\textit{APMEP / CoopMaths. CC BY-SA. Présentation adaptée par Maths Explorer.}\par\medskip
L'exercice est constitué de deux parties indépendantes.

\medskip

\noindent {\Large \textbf{Partie 1}}

\medskip

On considère l'équation différentielle :
\[(E): y' + y = \e^{-x}\]

\begin{enumerate}
\item Soit $u$ la fonction définie sur $\R$ par $u(x) = x \e^{-x}$.

Vérifier que la fonction $u$ est une solution de l'équation différentielle $(E)$.
\item On considère l'équation différentielle $\left(E'\right): y' + y = 0$.

Résoudre l'équation différentielle $\left(E'\right)$ sur $\R$.

\item En déduire toutes les solutions de l'équation différentielle $(E)$ sur $\R$.

\item Déterminer l'unique solution $g$ de l'équation différentielle $(E)$ telle que $g(0) = 2$.
\end{enumerate}


\medskip

\noindent {\Large \textbf{Partie 2}}

\medskip

Dans cette partie, $k$ est un nombre réel fixé que l'on cherche à déterminer.

On considère la fonction $f_{k}$ définie sur $\R$ par :
\[f_{k}(x) = (x + k) \e^{-x}.\]

Soit $h$ la fonction définie sur $\R$ par
\[h(x) = \e^{-x}.\]

On note $C_{k}$ la courbe représentative de la fonction $f_{k}$ dans un repère orthogonal et $C$ la courbe représentative de la fonction $h$.

On a représenté sur le graphique en fin d'exercice les courbes $C_{k}$ et $C$ sans indiquer les unités sur les axes ni le nom des courbes.

	\begin{enumerate}
		\item Sur le graphique en fin d'exercice à rendre avec la copie, l'une des courbes est en traits pointillés, l'autre est en trait plein. Laquelle est la courbe $C$ ?

		\item En expliquant la démarche utilisée, déterminer la valeur du nombre réel $k$ et placer sur le graphique ci-dessous avec la copie l'unité sur chacun des axes du graphique.
	\end{enumerate}

\begin{center}
	\begin{tikzpicture}
		\begin{axis}[x = {20mm},y = {20mm}, xmin = -2.5, xmax = 4.5, ymin = -0.5, ymax = 3,
			xtick = \empty, ytick = \empty,
			grid = major, axis lines = center]
			\addplot [line width = 1.2pt,color = red,smooth,samples = 100,domain = -2.5:4.5 ]{exp(-x)};
			\addplot [line width = 1.2pt,color = teal,dashed,smooth,samples = 100,domain = -2.5:4.5 ]{(x + 2)*exp(-x)};
		\end{axis}
	\end{tikzpicture}
\end{center}

\end{document}
